\chapter{术语表和设计不变量}

本章把开发者在维护 CCB 时反复遇到的术语和不变量集中说明。它的作用是减少跨模块沟通成本，避免同一个词在 CLI、daemon、provider 和 message bureau 中被不同理解。

\section{项目和运行时术语}

\begin{longtable}{p{0.25\linewidth}p{0.65\linewidth}}
\toprule
术语 & 含义 \\
\midrule
project anchor & 项目根目录下的 \src{.ccb}。它定义一个 CCB 项目的边界。 \\
backend & 项目范围的 \term{ccbd} daemon。一个 anchor 只有一个权威 backend。 \\
keeper & 项目级生命周期守护者，负责推进 backend generation 和监督启动。 \\
generation & backend 生命周期代次，用来区分当前 daemon 与旧 residue。 \\
mounted & 当前 generation 有 ready project socket，不能只看旧 socket path。 \\
configured agents & \src{ccb.config} 中期望 daemon 挂载和监督的 agent set。 \\
foreground layout & 用户看到的 tmux window/pane 布局。它是 UI，不是全部状态。 \\
runtime binding & agent 与 provider runtime、pane、session、workspace 的绑定事实。 \\
runtime health & daemon registry 对 provider runtime 当前健康状态的判断。 \\
residue & 旧 pane、旧 pid、旧 socket、未知 agent 目录等非权威残留证据。 \\
\bottomrule
\end{longtable}

\section{通信术语}

\begin{longtable}{p{0.25\linewidth}p{0.65\linewidth}}
\toprule
术语 & 含义 \\
\midrule
ask & 用户或 agent 提交给另一个 agent 的任务请求。 \\
MessageEnvelope & CLI 到 daemon 的请求边界对象。 \\
submission & 一次 logical submit，可包含多个 jobs。 \\
message & 通信层的逻辑请求，带 reply policy 和 retry policy。 \\
attempt & message 对某个 agent 的一次尝试，对应一个 job。 \\
job & dispatcher 执行单元。 \\
inbound event & mailbox 中的 task request 或 task reply 事件。 \\
reply & finalization 写入的通信结果。 \\
callback edge & parent job、child job 和 continuation 的持久关系。 \\
reply delivery & 把 reply event 自动投递给目标 agent 的特殊 job 流程。 \\
artifact & 长文本 payload 的存储间接层。 \\
\bottomrule
\end{longtable}

\section{状态不变量}

下面的不变量比单个实现细节更稳定。

\begin{enumerate}
  \item 一个 \src{.ccb} anchor 只拥有一个权威 backend。
  \item \src{ccb.config} 定义期望 agent mount set 和前台布局。
  \item backend lifecycle records 和 current generation socket readiness 决定 mounted。
  \item residue 不能阻止 bootstrap，也不能重新定义 project state。
  \item pane death 是 daemon supervision concern。
  \item \cmd{cmd} 是 foreground anchor pane，不是 provider agent。
  \item provider session files 是证据，不是唯一权威。
  \item dispatcher queue 和 mailbox queue 是不同队列。
  \item completion decision 是 job terminal 的依据。
  \item reply record 是通信层对 terminal decision 的投影。
  \item callback 是持久 edge，不是同步 wait。
  \item retry 延续原 message lineage，resubmit 创建新 lineage。
  \item ack 只能确认 inbox head。
  \item artifact 只是 payload indirection。
  \item source runtime validation 不能污染 source checkout 的 active CCB runtime。
\end{enumerate}

\section{配置不变量}

配置系统有自己的不变量。

\begin{itemize}
  \item config version 必须是 2。
  \item compact、rich、hybrid 最终都要进入同一个 validation 和 ProjectConfig。
  \item hybrid overlay 只能覆盖 compact layout 已有 agent 的非 header-owned 字段。
  \item legacy \src{default_agents/layout/cmd_enabled} 与 windows topology 互斥。
  \item windows leaf 必须声明 provider。
  \item windows topology 中不支持 \src{cmd} leaf。
  \item tool windows 要求 windows topology。
  \item role id shorthand 必须解析到已安装 role 的 default agent name。
  \item API shortcut 不能与 inherited auth/config 或显式 API env 冲突。
\end{itemize}

\section{provider 不变量}

\begin{itemize}
  \item provider 是 agent 的执行后端，不是 agent 本身。
  \item runtime home 隔离和继承策略必须在 provider profile 中表达。
  \item provider-specific completion detector 不能用通用字符串规则替代。
  \item provider resume support 必须显式进入 retry policy 判断。
  \item provider command template 是高级逃逸口，不能破坏 session binding 假设。
\end{itemize}

\section{诊断不变量}

\begin{itemize}
  \item doctor 输出应帮助定位层次，而不是只给单一 ok/failed。
  \item storage diagnostics 应区分 missing、error、stale 和 residue。
  \item queue/inbox degraded 不等于通信事实全部丢失。
  \item trace 是连接 job、message、attempt、reply、event 的标准诊断入口。
  \item recovery 输出应包含 noop reason 或 changed audit，不应只说 success。
\end{itemize}

\section{维护用语规范}

写 issue、commit message 或文档时，建议使用以下表达：

\begin{itemize}
  \item “backend not mounted”，表示当前 generation socket 未 ready。
  \item “runtime unhealthy”，表示 registry/provider runtime health 异常。
  \item “mailbox summary missing”，表示 summary 文件缺失，但不等于 message store 丢失。
  \item “attempt terminal but reply missing”，表示 finalization 到 message-bureau reply 投影之间断裂。
  \item “reply head blocked”，表示 caller mailbox head 无法消费或 delivery。
  \item “role shorthand expansion”，表示配置加载阶段把 role id 转为 agent spec。
\end{itemize}

不要把所有问题都写成 “agent crashed”。这个说法太粗，会掩盖 daemon、provider、dispatcher、mailbox、reply delivery 和 UI 之间的边界。

\section{开发者检查清单}

改动前：

\begin{enumerate}
  \item 识别改动属于 CLI、config、daemon、dispatcher、message bureau、provider、storage、UI 的哪一层。
  \item 找到对应设计契约。
  \item 找到最小测试组。
  \item 判断是否影响 runtime state 或只影响显示。
\end{enumerate}

改动中：

\begin{enumerate}
  \item 不要把残留证据提升为权威。
  \item 不要让 CLI service 直接绕过 daemon 修改项目状态。
  \item 不要把 provider-specific 行为写进通用 completion 逻辑。
  \item 不要跳过 message-bureau lineage。
\end{enumerate}

改动后：

\begin{enumerate}
  \item 更新契约文档或手册。
  \item 运行主题测试。
  \item 如果涉及 source runtime，使用外部测试项目和隔离 HOME。
  \item 记录新的证据路径。
\end{enumerate}
